\begin{figure}[htbp]
\centering
\begin{tikzpicture}[x=.58cm,y=.58cm]
 \foreach \etapa/\dx in {2/0,4/10.4}{
  \begin{scope}[shift={(\dx,0)}]
   \bcuadricula{6}
   \breina{6}{1}{2}
   \breina{6}{2}{4}
   \ifnum\etapa=2\relax
    \def\casillas{3/1,3/6,4/1,4/3,5/3,5/5,6/1,6/3,6/5,6/6}
    \def\corte{b=2/3}
   \else
    \breina{6}{3}{6}
    \breina{6}{4}{1}
    \def\casillas{5/3,5/5,6/5}
    \def\corte{b'=1/3}
   \fi
   \foreach \fila/\col in \casillas{
    \draw[btinta,line width=.55pt] (\col-.5,6.5-\fila) circle[radius=.16];
    \bmarca{6}{\fila}{\col}
   }
   \draw[btinta,line width=.95pt] (0,6-\etapa)--(6,6-\etapa);
   \node[font=\small,anchor=north] at (3,-.18) {$\corte,\quad m=\the\numexpr6-\etapa\relax$};
  \end{scope}
 }
 \draw[->,>=stealth,btinta,line width=.8pt] (6.55,3)--(9.45,3)
   node[midway,above,font=\scriptsize] {rows 3 and 4};
\end{tikzpicture}
\caption{Two cut-offs of the solution \((2,4,6,1,3,5)\). Entry \(r\) gives the column of the queen in row
\(r\). In this example the rows are revealed from top to bottom, and the horizontal line marks how far
the process has gone. Crowns indicate revealed queens, and circles with a central dot indicate the
squares those
queens leave available. After rows \(3\) and \(4\) are revealed, three of the ten remain.}
\label{fig:cortes}
\end{figure}
